% Part: normal-modal-logic
% Chapter: syntax-and-semantics
% Section: introduction

\documentclass[../../../include/open-logic-section]{subfiles}

\begin{document}

\olfileid[fr]{nml}{syn}{int}

\olsection{Introduction}

La logique modale étudie les \emph{propositions modales} et les
relations de conséquence qui les unissent. En voici quelques exemples :
\begin{enumerate}
\item Il est nécessaire que $2+2=4$.
\item Il est nécessairement possible qu'il pleuve demain.
\item Si~$!A$ est nécessairement possible, alors~$!A$ est possible.
\end{enumerate}
La possibilité et la nécessité ne sont pas les seules modalités :
d'autres connecteurs unaires sont également classés parmi elles,
par exemple « il est obligatoire que~$!A$ »,
« il sera vrai que~$!A$ »,
« Dana sait que~$!A$ » et « Dana croit que~$!A$ ».

La logique modale apparaît déjà dans le \emph{De Interpretatione}
d'Aristote. Celui-ci a été le premier à remarquer que la nécessité
implique la possibilité, mais non l'inverse; que possibilité et
nécessité se définissent l'une par l'autre; que, s'il est possible que
$!A \land !B$ soit vrai, alors il est possible que $!A$ soit vrai
et que $!B$ soit vrai, sans que la réciproque vaille; et que, si
$!A \to !B$ est nécessaire, la nécessité de~$!A$ entraîne celle
de~$!B$.

La première approche moderne de la logique modale est due à C.~I.
Lewis. Ses travaux culminent dans \emph{Symbolic Logic} (1932),
écrit avec Langford. Lewis et Langford ne se satisfaisaient pas de la
représentation de l'implication par le conditionnel matériel :
$!A \lif !B$ rend mal « $!A$ implique $!B$ ». Ils proposaient
plutôt de caractériser l'implication par « nécessairement, si $!A$,
alors $!B$ », ce que l'on note $!A \strictif !B$. En cherchant à
distinguer les propriétés de ces notions, Lewis a identifié cinq
systèmes modaux, \Log{S1}, \ldots, \Log{S4} et \Log{S5}, dont les
deux derniers sont encore utilisés.

L'approche de Lewis et Langford était purement \emph{syntaxique} :
ils choisissaient des axiomes et des règles raisonnables, puis
examinaient ce que ceux-ci permettaient de démontrer. Une approche
sémantique est longtemps restée difficile à établir. Rudolf Carnap
en a proposé une première dans \emph{Meaning and Necessity} (1947),
à l'aide de la notion de \emph{description d'état}, c'est-à-dire
d'un ensemble d'énoncés atomiques tenus pour « vrais » dans cet état.
Après avoir étendu la définition de la vérité aux énoncés
quelconques~$!A$, Carnap dit que~$!A$ est \emph{nécessairement vrai}
s'il est vrai dans toutes les descriptions d'état. Cette approche
ne permettait pas de distinguer les modalités \emph{itérées} :
les énoncés de la forme « Il est possible qu'il soit nécessaire que
\ldots{} il soit possible que~$!A$ » se ramenaient toujours à la
modalité la plus intérieure.

La sémantique modale connaît une avancée majeure avec le théorème de
complétude pour~$\Log{S5}$ que Saul Kripke publie en 1959 dans
« A Completeness Theorem in Modal Logic » (\emph{Journal of Symbolic Logic}).
La présentation par une relation d'accessibilité figure dans
\href{https://doi.org/10.1002/malq.19630090502}{\emph{Semantical Analysis of Modal Logic I} (1963), p.~68--70}.
Kripke partait de l'idée de Leibniz selon laquelle un énoncé est
nécessairement vrai s'il est vrai « dans tous les mondes possibles ».
Mais si l'on quantifie ainsi sur tous les mondes sans tenir compte
du monde~$w$ où l'on évalue l'énoncé (ou de la description
d'état~$s$ dans l'approche de Carnap), la vérité d'un énoncé précédé
de « nécessairement » ne dépend pas de~$w$. La vérité d'un atome ou
d'une formule où un atome se trouve hors de la portée d'une modalité
peut toutefois varier d'un monde à l'autre.

Kripke a donc supposé les mondes reliés par une \emph{relation d'accessibilité}~$R$ : un
énoncé de la forme « Nécessairement~$!A$ » est vrai en un monde~$w$
si et seulement si~$!A$ est vrai en tous les mondes~$w'$
\emph{accessibles depuis}~$w$. Les sémantiques fondées sur une
version de cette idée portent le nom de \emph{sémantiques de Kripke}
et ont permis le développement foisonnant des logiques modales.

Interprétée par la sémantique de Kripke, la logique modale nous
montre à quoi ressemblent les \emph{structures relationnelles}
« de l'intérieur ». Une structure relationnelle est simplement un
ensemble muni d'une relation binaire : la classe des élèves ordonnés
selon leur numéro de sécurité sociale en est un exemple. Les
structures relationnelles se rencontrent cependant dans des
domaines très divers. Outre les états du monde reliés par une
possibilité relative, on peut considérer les états de connaissance
d'un agent reliés par une possibilité épistémique, ou les états d'un
système dynamique reliés par leurs transitions. La logique modale
permet de modéliser tous ces cas : le premier relève de la logique
modale ordinaire, ou \emph{aléthique}; les autres conduisent
notamment aux logiques épistémique et dynamique.

Nous nous concentrerons sur un aspect particulier, appelé
\emph{théorie de la correspondance}. Kripke a notamment découvert
que de nombreuses propriétés de la relation d'accessibilité~$R$,
comme la transitivité ou la symétrie, peuvent être caractérisées
\emph{dans le langage modal} lui-même au moyen de « schémas
modaux » appropriés. Les logiciens disent, par exemple, que la
réflexivité de~$R$ « correspond » au schéma « Si nécessairement~$!A$,
alors~$!A$ ». Nous étudierons surtout la théorie de la
correspondance pour plusieurs systèmes modaux classiques, dont
\Log{S4} et \Log{S5}, obtenus en combinant les schémas
\Ax{D}, \Ax{T}, \Ax{B}, \Ax{4} et~\Ax{5}.

Dans un cadre de Kripke, la validité du schéma selon lequel la
nécessité implique la possibilité correspond à la sérialité de~$R$.
Chaque monde doit alors avoir au moins un successeur accessible.

\end{document}
